\section{Introduction}
As an object approaches a camera, its apparent size grows. For rigid frontal motion at approximately constant relative speed, the ratio of apparent size $s$ to expansion rate $\dot s$ gives time-to-contact (TTC), $T\simeq s/\dot s$. A recognizer can identify the object while leaving this rate undetermined; estimating when contact will occur requires recovering the change in scale. Rotation, occlusion, and camera motion complicate that measurement because they can alter the image independently of closing distance.

RGB frames provide texture and spatial context, while event cameras report asynchronous brightness changes along moving contours between frames \citep{eventsurvey}. To exploit both, a fusion method must account for the timing and support of each observation as well as the different representations learned from the two streams. We address this problem by estimating motion within each modality and combining quantities with shared meanings.

In PHASE-TTC, each branch supplies phase estimates together with timing and support information to a small residual network that adjusts an RGB anchor. The signed phase coordinate gives the branches a common description of approach and recession over a reference interval, allowing each to keep its own hidden representation. We preserve apparent expansion through common object coordinates and use temporal refiners to incorporate observations available before each query.

We evaluate PHASE-TTC as a complete TTC system, then examine its readout with matched controls. The private benchmark establishes system performance. To investigate the contribution of fusion, we hold the final upstream estimators and temporal histories fixed, compare simple combinations of their outputs, and train equally sized correction heads with or without event information. A separate local experiment asks whether learned routing adds value once apparent box expansion is available.

\subsection{Contributions}
We make three contributions:
\begin{enumerate}
\item \textbf{A causal phase-fusion architecture.} Independent motion estimators and temporal refiners supply a 20-field descriptor of phase, timing, and observation support to a 1,955-parameter residual head anchored on RGB. In the evaluated data most timing and support fields are constant, so the head acts mainly on ten phase values (Appendix~\ref{sec:physical-fusion-fields}). Table~\ref{tab:component-contracts} describes the nine-component system and its supervision.
\item \textbf{A benchmark result with matched readout controls.} The submitted system achieves \OfficialPhasePublishedOverallMiD{} Overall MiD on the archived GarlTTC Test leaderboard, against \OfficialGarlOverallMiD{} for the Garl-TTC reference. We complement this complete-system result with equal-capacity, equal-history head comparisons and simple fusion baselines on fixed final-model features. These controls do not show a benefit from event access or from the phase parameterization, so the benchmark gain is attributed to the complete system and not to the fusion interface (Sections~\ref{sec:official}--\ref{sec:design-analysis}).
\item \textbf{Geometric adaptation and runtime analysis.} We compare box expansion, fixed mixtures, and learned routing on the same queries to measure what each adds. A paced replay on a laptop GPU then measures numerical parity and end-to-end response against a 20\,ms deadline (Sections~\ref{sec:quality-adaptation}--\ref{sec:runtime}).
\end{enumerate}
The supplement retains the earlier predictive-learning, temporal, and physical-model investigations. They explain the design's development without serving as substitutes for the final-system comparisons.

\subsection{Relation to prior work}
\paragraph{Time-to-contact from apparent expansion.} The ratio of image size to its rate of change gives TTC without metric depth or speed \citep{lee1976}, and direct methods recover it from brightness derivatives for a planar surface \citep{horn2007}. Learned frame-based estimators regress or classify motion in depth between two images: optical expansion upgrades optical flow to scale change \citep{yang2020expansion}, Binary TTC decides whether contact occurs within a set of time thresholds \citep{badki2021bittc}, and scale-matching correlation searches explicitly across scales \citep{ling2023tpcv,ling2024scaleflowpp}. TSTTC provides a large frame-only driving benchmark for object-level TTC \citep{shi2023tsttc}. These methods use two frames; we do not evaluate against them, and their scale-aware matching is absent from our producers.

\paragraph{Event-based TTC and motion.} Event cameras support TTC through asynchronous flow \citep{clady2014}, contrast maximization with scaling or divergence models \citep{gallego2018cmax,mcleod2022}, and geometric fitting with spatio-temporal registration \citep{eventaided}. Dense event representations such as voxel grids and learned spike tensors feed convolutional backbones \citep{zhu2019voxel,gehrig2019est}, including correlation-based optical flow \citep{gehrig2021eraft}; event--frame fusion has reduced latency in automotive detection \citep{gehrig2024dagr}. EvTTC provides synchronized observations and a dedicated TTC benchmark \citep{evttc}; MVSEC and DSEC illustrate event--frame acquisition in driving conditions \citep{zhu2018mvsec,dsec}. The eAP study introduces geometry-aware representations, the logarithmic motion-in-depth metric, and early and late event--RGB fusion in the Garl models \citep{eap}. \paragraph{Position of this work.} Our contribution is the particular estimator interface: named motion estimates, timing, and support drive a residual correction around RGB. The phase coordinate itself comes from the benchmark geometry. Relative to feature-level late fusion, the interface exposes the meanings of its inputs and allows separate producer fitting. The matched tests below measure the utility of this readout under fixed producers; a matched hidden-feature fusion comparison is still needed to establish superiority over that alternative.

I-JEPA predicts target representations from context \citep{ijepa}, while VICReg regularizes representation variance and covariance \citep{vicreg}. These methods motivate our predictive-learning controls. Work on slow features shows that predictable distractors can dominate a learned state \citep{slowfeatures}; our experiments test how nuisance dynamics affect the information available for TTC. They support conclusions about the tested objectives and settings rather than a general impossibility of predictive learning.
